\section{Mathematical Logic}
\begin{questions}
    % Q1
    \item \begin{partquestions}{\alph*}
        \item Not a statement; it is neither true nor false, it is just a sentence.
        \item A statement; it is definitely false.
        \item A statement; it can (eventually) be shown to be true or false.
        \item Not a statement; the truth or falsity of the statement depends on the value of $n$.
        \item A statement; it is definitely true.
        \item A statement; we can find such a number $n$, for instance $n = 3$, so we can confirm that this is true.
        \item Not a statement. This is tricky since, algebraically, it is true for all $x$. However, the fact remains that we can only evaluate the truth of this statement by substituting values of $x$ since $x$ is an unknown.
    \end{partquestions}

    % Q2
    \item We work from the inner-most bracket outwards. We note $p$ is true, $q$ is false, and $r$ is true.
    \begin{itemize}
        \item $p \lor q$ can be understood to be ``1 is a positive number \textbf{or} 1 is an even number''. This is true since 1 is indeed a positive number.
        \item $(p \lor q) \land r$ is ``(1 is a positive number or 1 is an even number) \textbf{and} 1 is an odd number''. Since $p \lor q$ is true and 1 is, indeed, an odd number, thus this statement is true.
        \item $\lnot((p \lor q) \land r)$ is false since $(p \lor q) \land r$ is true.
    \end{itemize}
    Hence, the overall statement ``$\lnot((p \lor q) \land r)$ is false'' is true.

    % Q3
    \item \begin{partquestions}{\roman*}
        \item $s$ is the statement ``20 is a multiple of 5''.
        \item The atomic statements are ``20 ends in a 0'' and ``20 ends in a 5''.
        \item $t$ is $(p \lor q) \implies s$.
        \item ``20 is a multiple of 5 if and only if 20 ends in a 0 or 5'', which is true.
    \end{partquestions}

    \pagebreak
    
    % Q4
    \item Again we consider a truth table.
    \begin{table}[H]
        \centering
        \begin{tabular}{|c|c||c|c||c|c|}
            \hline
            $\boldsymbol{p}$ & $\boldsymbol{q}$ & $\boldsymbol{\lnot p}$ & $\boldsymbol{\lnot q}$ & $\boldsymbol{p \implies q}$ & $\boldsymbol{(\lnot q) \implies (\lnot p)}$ \\ \hhline{|=|=#=|=#=|=|}
            T & T & F & F & T & T \\ \hline
            T & F & F & T & F & F \\ \hline
            F & T & T & F & T & T \\ \hline
            F & F & T & T & T & T \\ \hline
        \end{tabular}
    \end{table}
    Therefore we establish the required result.

    % Q5
    \item The first statement is proven by considering the truth table below.
    \begin{table}[H]
        \centering
        \begin{tabular}{|c|c||c|c|c||c|c|}
            \hline
            $\boldsymbol{p}$ & $\boldsymbol{q}$ & $\boldsymbol{\lnot p}$ & $\boldsymbol{\lnot q}$ & $\boldsymbol{p \land q}$ & $\boldsymbol{\lnot(p \land q)}$ & $\boldsymbol{(\lnot p) \lor (\lnot q)}$ \\ \hhline{|=|=#=|=|=#=|=|}
            T & T & F & F & T & F & F \\ \hline
            T & F & F & T & F & T & T \\ \hline
            F & T & T & F & F & T & T \\ \hline
            F & F & T & T & F & T & T \\ \hline
        \end{tabular}
    \end{table}
    The second statement is proven by considering the truth table below.
    \begin{table}[H]
        \centering
        \begin{tabular}{|c|c||c|c|c||c|c|}
            \hline
            $\boldsymbol{p}$ & $\boldsymbol{q}$ & $\boldsymbol{\lnot p}$ & $\boldsymbol{\lnot q}$ & $\boldsymbol{p \lor q}$ & $\boldsymbol{\lnot(p \lor q)}$ & $\boldsymbol{(\lnot p) \land (\lnot q)}$ \\ \hhline{|=|=#=|=|=#=|=|}
            T & T & F & F & T & F & F \\ \hline
            T & F & F & T & T & F & F \\ \hline
            F & T & T & F & T & F & F \\ \hline
            F & F & T & T & F & T & T \\ \hline
        \end{tabular}
    \end{table}

    % Q6
    \item One sees that
    \begin{align*}
        &((p \lor \lnot q) \land \lnot r) \lor p\\
        &\equiv p \lor ((p \lor \lnot q) \land \lnot r) & (\text{Commutativity})\\
        &\equiv (p \lor (p \lor \lnot q)) \land (p \lor \lnot r) & (\text{Distributivity})\\
        &\equiv ((p \lor p) \lor \lnot q) \land (p \lor \lnot r) & (\text{Associativity})\\
        &\equiv (p \lor \lnot q) \land (p \lor \lnot r) & (\text{Clearly } p \lor p \equiv p)\\
        &\equiv p \lor (\lnot q \land \lnot r) & (\text{Distributivity})\\
        &\equiv p \lor \lnot(q \lor r). & (\text{De Morgan's Law})
    \end{align*}

    % Q7
    \item One formulation is ``$\forall n, ((n > 2) \implies (\forall x, y, z, (x^n + y^n \neq z^n)))$''.
    
    % Q8
    \item We use a truth table.
    \begin{table}[H]
        \centering
        \begin{tabular}{|c|c||c|c||c|c|}
            \hline
            $\boldsymbol{p}$ & $\boldsymbol{q}$ & $\boldsymbol{\lnot q}$ & $\boldsymbol{p \implies q}$ & $\boldsymbol{\lnot(p \implies q)}$ & $\boldsymbol{p \land \lnot q}$ \\ \hhline{|=|=#=|=#=|=|}
            T & T & F & T & F & F \\ \hline
            T & F & T & F & T & T \\ \hline
            F & T & F & T & F & F \\ \hline
            F & F & T & T & F & F \\ \hline
        \end{tabular}
    \end{table}

    % Q9
    \item \begin{partquestions}{\roman*}
        \item One formulation is ``$\forall x, ((x \neq 0) \implies (\exists y \text{ s.t. } xy = 1))$''.
        \item The negation of the above formulation is ``$\exists x \text{ s.t. } ((x \neq 0) \land (\forall y, xy \neq 1))$''.
    \end{partquestions}
\end{questions}
